\documentclass[11pt,a4paper]{article}
\usepackage{graphicx}
\usepackage{xcolor}
\usepackage{geometry}
\usepackage{amsmath}
\usepackage{booktabs}
\usepackage{tabularx}
\usepackage{adjustbox}
\usepackage{enumitem}
\usepackage{caption}
\usepackage{microtype}
\usepackage{xurl}
\usepackage{placeins}

\geometry{a4paper, top=2.6cm, bottom=2.6cm, left=2.5cm, right=2.5cm}
\captionsetup{font=small, labelfont=bf, skip=6pt}
\widowpenalty=10000
\clubpenalty=10000
\emergencystretch=1.8em
\raggedbottom


\usepackage[
    colorlinks=true,
    linkcolor=blue!55!black,
    citecolor=blue!45!black,
    urlcolor=blue!55!black,
    bookmarks=true,
    unicode=true,
    pdftitle={Companion Review: Research Assessment, Code Audit, and Recommendations},
    pdfauthor={Prepared as companion to the SF-9 research paper},
    pdfsubject={Code audit and research assessment}
]{hyperref}

\newcommand{\code}[1]{\texttt{#1}}

\begin{document}

\begin{center}
    {\color{black!85}\rule{\textwidth}{0.4pt}}\\[1.6em]
    {\LARGE\bfseries Companion Review\par}
    \vspace{0.4em}
    {\large Research Assessment, Code Audit, and Recommendations\par}
    \vspace{0.8em}
    {\small for the repository \code{Daveshvats/}\\
    \code{Federated-Learning-for-Distributed-Space-Weather-Monitoring} (SF-9)\par}
    \vspace{1.0em}
    {\color{black!85}\rule{\textwidth}{0.4pt}}
\end{center}

\vspace{0.6em}

\section*{1. Overall Rating: 7.0 / 10 --- ``Promising first application with a
clear improvement path''}

\begin{table}[h]
\centering
\small
\begin{tabularx}{0.94\textwidth}{@{}l c X@{}}
\toprule
\textbf{Dimension} & \textbf{Score} & \textbf{Rationale} \\
\midrule
Problem importance & 9/10 & GIC risk to grids is real, documented
(1989 Qu\'ebec; 2012 near-miss), and the data-sovereignty barrier is
genuine. \\
Novelty & 8/10 & First FL treatment of solar flare prediction; verified
against the current literature during paper preparation. Well-targeted
cross-silo scenario. \\
Methodology & 6.5/10 & Sound framework (Dirichlet non-IID, FedProx,
Fed-Focal, $F_\beta$ thresholds), but the evaluation carries
methodological debts: single seed, thresholds tuned on the test set,
non-disjoint shards, and one stale metric. \\
Results strength & 5.5/10 & The direction is right (FedProx AUC 0.930
vs centralised 0.978) but absolute skill (F1 0.326, recall 0.416) is
sentinel-grade, not forecaster-grade; FedAvg's collapse is instructive
but weakens the headline. \\
Reproducibility & 7.5/10 & Deterministic seed, versioned fixes
(v2.1--v2.6), committed result figures; but no README, placeholder
values in \code{ABSTRACT.md}, and metrics exist only as PNG images. \\
Code quality & 7/10 & Exceptional inline documentation and diagnostic
commentary; offset by the audit findings in \S3. \\
\midrule
\textbf{Overall} & \textbf{7.0/10} & A credible workshop/short-conference
contribution; the fixes in \S4 would lift it toward a strong
full-conference or journal paper. \\
\bottomrule
\end{tabularx}
\end{table}

\section*{2. What the research does well}

\begin{enumerate}[leftmargin=2em, itemsep=3pt]
  \item \textbf{The right regime.} Cross-silo FL with six institution-like
  clients is the correct abstraction for space-weather agencies; the
  privacy constraint is enforced by construction, not just claimed.
  \item \textbf{Honest engineering narrative.} The v2.1--v2.6 fix log
  (SCAFFOLD NaN root cause, focal-alpha drift, feature-schema audit) is
  rare, valuable documentation; the companion paper turns it into a
  transferable lessons section.
  \item \textbf{Safety-critical framing done correctly.} Recall-weighted
  $F_\beta$ ($\beta=2$) threshold optimisation, rather than default
  0.5-threshold metrics, matches how a grid operator would actually
  consume the model.
  \item \textbf{Discrimination retained.} FedProx holds 95\% of the
  centralised AUC (0.930/0.978) with zero raw-data movement --- the core
  feasibility result stands.
  \item \textbf{Interpretable, physical features.} Decoded SHAP ranks
  (R-value extrema, TOTUSJH, Lorentz-force variability) align with
  established flare physics, which strengthens scientific credibility.
\end{enumerate}

\section*{3. Code Audit Findings (cross-check of the released scripts)}

Ordered by severity; all were cross-checked against the committed
artefacts and the extracted result values.

\begin{enumerate}[leftmargin=2em, itemsep=4pt]
  \item \textbf{[Bug, medium] Stale accuracy for federated models.}
  In \code{main.py}, after $F_\beta$ threshold optimisation, the
  \code{accuracy} field of FedAvg/FedProx results is never recomputed
  (only the centralised models get
  \code{accuracy\_score(...)} on new predictions). The committed
  \code{Comparison\_Table.png} therefore reports FedProx accuracy as
  0.019 --- the value at the pre-optimisation threshold --- while its
  confusion matrix implies $\approx$0.969. The companion paper reports
  the corrected value with a footnote. \emph{Fix:} recompute accuracy
  (and store predictions) for FL models exactly as done for baselines.
  \item \textbf{[Bug, medium] Dirichlet shards are not disjoint.}
  The Dirichlet partitioner in \code{partition\_clients.py} draws, for
  every (client, class) pair, a fresh random subset of the class pool
  with \code{np.random.choice}; nothing excludes indices already
  assigned to other clients, so training samples can appear in
  multiple clients (and some samples are never used). This inflates
  effective duplication and weakens the ``six disjoint observatories''
  semantics. \emph{Fix:} draw a single permutation per class and slice
  it by the Dirichlet proportions.
  \item \textbf{[Bug, low-medium] Interpretability names lost.}
  The 144 statistical features are renamed \code{feature\_0..143} in
  \code{\_arrays\_to\_dataframe}, so the SHAP plot (and every downstream
  table) shows opaque indices; the companion paper had to decode them
  analytically. \emph{Fix:} propagate
  \code{mean\_R\_VALUE}-style names end to end.
  \item \textbf{[Transparency] ``FedAvg'' is not vanilla FedAvg.} Both
  \code{run\_fedavg} and \code{run\_fedprox} aggregate through
  \code{aggregate\_weights\_dafl} (distribution-aware weighting). The
  comparison therefore isolates the proximal term \emph{given} DA-FL
  weighting --- a legitimate design, but the labels in the committed
  figures understate this. \emph{Fix:} either add a true-FedAvg
  (size-weighted only) arm, or describe DA-FL in the figure caption.
  \item \textbf{[Config mismatch] LSTM vs committed outputs.}
  \code{config.py} ships \code{USE\_LSTM=True}, but every committed
  artefact (\code{FedAvg MLP}, \code{FedProx MLP}) is from the MLP path
  with 144 statistical features. Anyone re-running ``as released'' will
  get a different pipeline than the one that produced the figures.
  \emph{Fix:} pin the configuration that reproduces the committed
  outputs (or commit the LSTM outputs too).
  \item \textbf{[Methodological debt] Threshold selected on the test
  set.} The $F_\beta$ sweep of \code{main.py} runs directly on the test
  partitions. The optimism is bounded (one scalar from an 81-point
  grid) but reviewers will flag it. \emph{Fix:} select $\tau^\ast$ on a
  carved-out validation split, report final numbers on test only.
  \item \textbf{[Consistency, low]} \code{config.py}
  \code{BATCH\_SIZE=256} is unused by the FL loop
  (\code{\_make\_loader} hardcodes 512); \code{FOCAL\_ALPHA} in config
  is bypassed by the hardcoded \code{focal\_alpha=0.25} argument;
  \code{THRESHOLD=0.35} silently drives the misleading training-time
  convergence curves (the paper documents why these are uninformative).
  Consolidate these into one source of truth.
  \item \textbf{[Hygiene, low]} \code{ABSTRACT.md} still contains
  \code{[X.XX]} placeholder metrics; the repo has no README (setup,
  data download, and run instructions live only in docstrings); results
  exist only as PNGs, not JSON/CSV, which blocks automated
  re-comparison.
\end{enumerate}

\section*{4. Recommendations (prioritised)}

\begin{description}[leftmargin=2em, itemsep=4pt]
  \item[P0 --- before any submission:] fix findings 1--2 above; add
  \code{F2} (the optimised metric!) to the reported table; run the
  five-seed repetition with means $\pm$ std; regenerate the figures from
  the fixed code.
  \item[P1 --- strengthens the science:] add a \emph{centralised MLP}
  baseline (currently the federated model is compared only against a
  300-tree XGBoost, which conflates architecture with privacy cost);
  evaluate the LSTM path; re-enable SCAFFOLD now that its SGD
  control-variate fix is in; replace the threshold sweep with prior
  correction / logit adjustment so calibration is principled rather
  than post hoc.
  \item[P2 --- for the venue:] a $\mu$-sweep ($10^{-3}$--$10^{-1}$) and
  $\alpha$-sweep (0.5--5) to map the stability boundary; per-client
  evaluation (is the global model better than each client's local
  model?); HARP-based geographic partitioning for realism.
  \item[P3 --- deployment track:] secure aggregation and differential
  privacy on the update channel (the current design protects data
  locality, not gradient-inversion attacks); Ditto-style personalised
  heads; SCADA integration notes.
\end{description}

\section*{5. Verdict}

The repository delivers a genuinely novel, well-motivated, and
honestly-engineered first study of federated solar flare prediction,
with a usable feasibility result (95\% of centralised AUC without data
movement) and unusually good in-code documentation. Its weaknesses ---
single-seed evidence, sentinel-grade absolute skill, several
audit-level bugs, and metrics recoverable only from images --- are all
fixable with bounded effort, and none undermines the core claim. With
the P0/P1 items addressed, this is a solid conference paper; the
companion manuscript is written to that standard so the numbers it
reports will survive the fixes (the one corrected value is footnoted
at its source).

\end{document}
